%=======================================================================
%  IE 5311-001 (D01) Exploration Problem EP8
%  One page, single spacing, 10 pt, 1-inch margins, as the sheet requires.
%  Content past page 1 is not graded, so the page limit is hard.
%  Sections map onto the three parts asked: the general formulation,
%  two reformulations, and their relationship / merits / limits.
%
%  IF IT SPILLS ONTO A SECOND PAGE, delete in this order until it fits.
%    1. the last sentence of section 5 (from "For nonlinear")
%    2. the K = 12 row of the table in section 4
%    3. the parenthesis in section 1 beginning "(Raman and Grossmann"
%  Do not cut section 3 or section 4: the relationship and the measured
%  gap are what the question asks for.
%=======================================================================
\documentclass[10pt]{article}

\usepackage[margin=1in]{geometry}
\usepackage{amsmath}
\usepackage{amssymb}
\usepackage{booktabs}
\usepackage{array}

\DeclareMathOperator{\conv}{conv}
\newcommand{\T}{^{\!\top}}

\pagestyle{empty}
\setlength{\parskip}{0.35em}
\setlength{\parindent}{0pt}

\begin{document}

{\large\bfseries EP8.\quad Disjunctive programming and its two reformulations}\\[1pt]
{\small Lidivine Kengne \quad\textbar\quad IE 5311-001 (D01), Homework 6}
\vspace{0.3em}

\textbf{1. The general formulation.}
A disjunctive program optimises over a \emph{union} of polyhedra rather
than one. With $K$ disjuncts,
\[
  \min_{x}\; c\T x
  \qquad\text{s.t.}\qquad
  x \in \bigcup_{k \in K} P_k, \quad P_k = \{x : A_k x \le b_k\},
  \qquad x \in X \subseteq \mathbb{R}^n,
\]
$X$ a bounding box. The union is the whole difficulty: each $P_k$ is
convex, their union is not. Every complementarity condition is a
two-term instance, since $ab = 0$ with $a,b \ge 0$ is
$\{a = 0\} \cup \{b = 0\}$ --- which is why Problem~2 of this homework
is a disjunctive program in disguise.

\textbf{2. The two reformulations.}
Both introduce binaries $z_k$ with $\sum_k z_k = 1$ and differ in how
they switch the constraints off.

\emph{Big-M.} Keep one copy of $x$ and relax each inactive block by a
constant: $A_k x \le b_k + M_k(1 - z_k)$. The disjunct is enforced when
$z_k = 1$ and vacuous when $z_k = 0$, provided $M_k$ is large enough
that $A_k x \le b_k + M_k$ holds throughout $X$.

\emph{Hull, or disaggregated, reformulation.} Give each disjunct its
own copy of the variable and split $x$ among them:
\[
  x = \sum_{k \in K} x^k, \qquad
  A_k x^k \le b_k z_k, \qquad
  \sum_{k \in K} z_k = 1, \qquad z \ge 0 .
\]
A disjunct with $z_k = 0$ forces $x^k = 0$.

\textbf{3. The relationship.}
Balas' theorem: the projection of the hull relaxation onto $x$ is
exactly $\conv\bigl(\bigcup_k P_k\bigr)$, the tightest convex
relaxation there is. The big-M relaxation always contains it, so its
bound is never better, and how much worse depends entirely on how close
each $M_k$ is to the smallest valid value.

There is an instructive exception. For the two-term complementarity
with $a \le U_a$, $b \le U_b$, eliminating $z$ from $a \le U_a z$,
$b \le U_b(1-z)$, $z \in [0,1]$ leaves $a/U_a + b/U_b \le 1$: the
triangle on $(0,0),(0,U_b),(U_a,0)$, which \emph{is} the hull of the
two segments. With tight bounds big-M already is the hull formulation.
Inflating $M$ destroys that:

\begin{center}\footnotesize
\setlength{\tabcolsep}{8pt}
\begin{tabular}{@{}rrrr@{}}
\toprule
$M_a, M_b$ relative to the tight bound & relaxation area
  & $/\,$hull & $/\,$variable box\\
\midrule
$1.00\times$ & 16.00 & $1.00\times$ & 0.50\\
$1.25\times$ & 23.00 & $1.44\times$ & 0.72\\
$2.00\times$ & 32.00 & $2.00\times$ & 1.00\\[-1pt]
\bottomrule
\end{tabular}
\end{center}

At twice the tight bound the relaxation is the whole variable box: the
complementarity has stopped constraining it. That is the precise reason
Problem~2 Part~3 had to \emph{derive} $M^{\mathrm{dual}} = 8$ rather
than pick a round number.

\textbf{4. Where they genuinely differ.}
Minimising $c\T x$ over a union of $K$ disjoint boxes in
$\mathbb{R}^2$, presolve and cuts off so the nodes mean something
(\texttt{ep8\_disjunctive.py}):

\begin{center}\footnotesize
\setlength{\tabcolsep}{6pt}
\begin{tabular}{@{}rrrrrrr@{}}
\toprule
$K$ & optimum & big-M LP & hull LP & big-M gap
  & big-M nodes & hull nodes\\
\midrule
 3 & 10.2208 & 8.0908 & 10.2208 & 2.1300 &  3 & 1\\
12 &  3.4317 & 2.1526 &  3.4317 & 1.2792 & 15 & 1\\
20 &  3.3480 & 0.4090 &  3.3480 & 2.9390 & 22 & 1\\
\bottomrule
\end{tabular}
\end{center}

The hull relaxation returns the integer optimum in every case and
closes at the root; big-M leaves a gap that branching must close, and
its node count grows with $K$ while the hull's does not.

\textbf{5. Advantages and limitations.}
Big-M is compact --- $|K|$ binaries, no new continuous variables ---
and portable to any MILP solver. Against that it needs valid bounds,
which can be costly or impossible to compute tightly: too large is
valid but useless, too small silently removes feasible points, and
that is the dangerous failure because the model still solves and still
returns a number. The hull needs no bounds and gives the best possible
relaxation, but pays $n|K|$ extra continuous variables, so it wins
only when the tighter bound removes more nodes than the larger LP
costs. For nonlinear disjuncts it also needs the perspective function
$z f(x/z)$, nonsmooth at $z = 0$.

\end{document}
